% GENERATED FILE. Edit canonical lessons, then run npm run generate:books.
% canonical-source-hash: fnv1a64:e1da58cec4c5dcf8
% canonical-lessons: ES-C700-nevar, ES-C700-helar, ES-C700-ventisca, ES-C700-nevada, ES-C700-escarcha

\chapter{Nieve y hielo}
\label{ch:es-nieve-y-hielo}
\hlchaptermodality{700}

\begin{chapteropening}
\textbf{By the end of this chapter:} \emph{I can talk about winter weather: snow, snowfalls, blizzards and freezing.}

Complete the last lesson of chapter 700: I can talk about winter weather: snow, snowfalls, blizzards and freezing.
\end{chapteropening}

\section[nevar]{nevar — to snow}

\label{lesson:ES-C700-nevar}



\begin{quote}
Before the new one: say the Spanish for to hail, then the Spanish for to thunder.
\end{quote}

\subsection*{You'll want to know: nevar}
\textbf{nevar} — \textquotedblleft{}to snow\textquotedblright{}.

\begin{quote}
En la sierra nieva en enero.
\end{quote}

\begin{quote}
\emph{It snows in the mountains in January.}
\end{quote}

The e becomes ie: \emph{nieva}.

\begin{cousinweb}
\emph{nevar} is inherited from Vulgar Latin \emph{nivāre} \textquotedblleft{}to snow\textquotedblright{} (EtymDB), from Latin \emph{nix, nivis} \textquotedblleft{}snow\textquotedblright{}.
\end{cousinweb}

\begin{grammarlens}[title={What you've built}]
One more word for telling, explaining and describing at B1.
\end{grammarlens}

\subsection*{Your turn}
\begin{itemize}
\raggedright
  \item \emph{Say it:} \emph{nevar}
  \item \emph{Say it:} \emph{nevar}, once more
  \item \emph{Say it:} a sentence of your own with \emph{nevar}
  \item \emph{Recall:} say the Spanish for to hail, then the Spanish for to thunder, then say \emph{nevar} again
\end{itemize}

\begin{tcolorbox}[breakable,colback=teal!4,colframe=teal!35!black,title={Before you move on}]
What does nevar mean? (\textquotedblleft{}To snow\textquotedblright{}.) Say it once more.
\end{tcolorbox}
\section[helar]{helar — to freeze}

\label{lesson:ES-C700-helar}



\begin{quote}
Before the new one: say the Spanish for to thunder, then the Spanish for to snow.
\end{quote}

\subsection*{You'll want to know: helar}
\textbf{helar} — \textquotedblleft{}to freeze\textquotedblright{}.

\begin{quote}
Anoche heló y la calle está blanca.
\end{quote}

\begin{quote}
\emph{It froze last night and the street is white.}
\end{quote}

The e becomes ie: \emph{hiela}. \emph{El helado} is ice cream.

\begin{cousinweb}
\emph{helar} is inherited from Latin \emph{gelāre} \textquotedblleft{}to freeze\textquotedblright{} (EtymDB), from \emph{gelū} \textquotedblleft{}frost\textquotedblright{}; Latin \emph{g} before \emph{e} became Spanish \emph{h} here.
\end{cousinweb}

\begin{grammarlens}[title={What you've built}]
One more word for telling, explaining and describing at B1.
\end{grammarlens}

\subsection*{Your turn}
\begin{itemize}
\raggedright
  \item \emph{Say it:} \emph{helar}
  \item \emph{Say it:} \emph{helar}, once more
  \item \emph{Say it:} a sentence of your own with \emph{helar}
  \item \emph{Recall:} say the Spanish for to thunder, then the Spanish for to snow, then say \emph{helar} again
\end{itemize}

\begin{tcolorbox}[breakable,colback=teal!4,colframe=teal!35!black,title={Before you move on}]
What does helar mean? (\textquotedblleft{}To freeze\textquotedblright{}.) Say it once more.
\end{tcolorbox}
\section[la ventisca]{la ventisca — the blizzard, the snowstorm}

\label{lesson:ES-C700-ventisca}



\begin{quote}
Before the new one: say the Spanish for to snow, then the Spanish for to freeze.
\end{quote}

\subsection*{You'll want to know: la ventisca}
\textbf{la ventisca} — \textquotedblleft{}the blizzard, the snowstorm\textquotedblright{}.

\begin{quote}
Con la ventisca no se ve nada.
\end{quote}

\begin{quote}
\emph{In the blizzard you cannot see a thing.}
\end{quote}

Snow and strong wind together.

\begin{cousinweb}
\emph{ventisca} is built on \emph{viento} \textquotedblleft{}wind\textquotedblright{}, which is inherited from Latin \emph{ventus}.
\end{cousinweb}

\begin{grammarlens}[title={What you've built}]
One more word for telling, explaining and describing at B1.
\end{grammarlens}

\subsection*{Your turn}
\begin{itemize}
\raggedright
  \item \emph{Say it:} \emph{la ventisca}
  \item \emph{Say it:} \emph{la ventisca}, once more
  \item \emph{Say it:} a sentence of your own with \emph{la ventisca}
  \item \emph{Recall:} say the Spanish for to snow, then the Spanish for to freeze, then say \emph{la ventisca} again
\end{itemize}

\begin{tcolorbox}[breakable,colback=teal!4,colframe=teal!35!black,title={Before you move on}]
What does la ventisca mean? (\textquotedblleft{}The blizzard, the snowstorm\textquotedblright{}.) Say it once more.
\end{tcolorbox}
\section[la nevada]{la nevada — the snowfall}

\label{lesson:ES-C700-nevada}



\begin{quote}
Before the new one: say the Spanish for to freeze, then the Spanish for the blizzard.
\end{quote}

\subsection*{You'll want to know: la nevada}
\textbf{la nevada} — \textquotedblleft{}the snowfall\textquotedblright{}.

\begin{quote}
Después de la nevada no había autobuses.
\end{quote}

\begin{quote}
\emph{After the snowfall there were no buses.}
\end{quote}

\emph{Una gran nevada} — heavy snow.

\begin{cousinweb}
\emph{nevada} is built on \emph{nevar} \textquotedblleft{}to snow\textquotedblright{}, which is inherited from Vulgar Latin \emph{nivāre}, from \emph{nix, nivis} \textquotedblleft{}snow\textquotedblright{}.
\end{cousinweb}

\begin{grammarlens}[title={What you've built}]
One more word for telling, explaining and describing at B1.
\end{grammarlens}

\subsection*{Your turn}
\begin{itemize}
\raggedright
  \item \emph{Say it:} \emph{la nevada}
  \item \emph{Say it:} \emph{la nevada}, once more
  \item \emph{Say it:} a sentence of your own with \emph{la nevada}
  \item \emph{Recall:} say the Spanish for to freeze, then the Spanish for the blizzard, then say \emph{la nevada} again
\end{itemize}

\begin{tcolorbox}[breakable,colback=teal!4,colframe=teal!35!black,title={Before you move on}]
What does la nevada mean? (\textquotedblleft{}The snowfall\textquotedblright{}.) Say it once more.
\end{tcolorbox}
\section[la escarcha]{la escarcha — the hoarfrost, the white frost}

\label{lesson:ES-C700-escarcha}



\begin{quote}
Before the new one: say the Spanish for the blizzard, then the Spanish for the snowfall.
\end{quote}

\subsection*{You'll want to know: la escarcha}
\textbf{la escarcha} — \textquotedblleft{}the hoarfrost, the white frost\textquotedblright{}.

\begin{quote}
Por la mañana hay escarcha en los coches.
\end{quote}

\begin{quote}
\emph{In the morning there is frost on the cars.}
\end{quote}

The thin white ice on grass and glass.

\begin{cousinweb}
The origin of \emph{escarcha} is not settled.
\end{cousinweb}

\begin{grammarlens}[title={What you've built}]
One more word for telling, explaining and describing at B1.
\end{grammarlens}

\subsection*{Your turn}
\begin{itemize}
\raggedright
  \item \emph{Say it:} \emph{la escarcha}
  \item \emph{Say it:} \emph{la escarcha}, once more
  \item \emph{Say it:} a sentence of your own with \emph{la escarcha}
  \item \emph{Recall:} say the Spanish for the blizzard, then the Spanish for the snowfall, then say \emph{la escarcha} again
\end{itemize}

\begin{tcolorbox}[breakable,colback=teal!4,colframe=teal!35!black,title={Before you move on}]
What does la escarcha mean? (\textquotedblleft{}The hoarfrost, the white frost\textquotedblright{}.) Say it once more.
\end{tcolorbox}
